\section{Algoritmo de muestreo DDPM}

\subsection{Muestreo del proceso inverso}

Una vez completada la fase de entrenamiento, el proceso para generar nuevos muestras (proceso inverso / proceso de muestreo) es el siguiente:

\begin{enumerate}
    \item Muestrear ruido inicial desde una distribución gaussiana estándar: $x_T \sim \mathcal{N}(0, I)$
    \item Para $t = T, T-1, \dots, 1$:
    \begin{enumerate}
        \item Predecir el ruido con la red: $\hat{\epsilon} = \epsilon_\theta(x_t, t)$
        \item Calcular la media: $\mu_\theta(x_t, t) = \frac{1}{\sqrt{\alpha_t}}\left(x_t - \frac{\beta_t}{\sqrt{1-\bar{\alpha}_t}}\, \hat{\epsilon}\right)$
        \item Si $t > 1$, muestrear $z \sim \mathcal{N}(0, I)$; si $t = 1$, establecer $z = 0$
        \item $x_{t-1} = \mu_\theta(x_t, t) + \sigma_t\, z$
    \end{enumerate}
    \item Retornar $x_0$
\end{enumerate}

\begin{importantbox}{Intuición del muestreo: eliminación progresiva de ruido}
El proceso de muestreo puede entenderse como: comenzar con un bloque de puro ruido y, en cada paso, permitir que la red «mire» la imagen actual para predecir cuánto ruido contiene, luego restar el ruido predicho (agregando una pequeña cantidad de nuevo ruido al mismo tiempo para mantener la aleatoriedad), iterando así hasta obtener una imagen limpia. El último paso ($t=1$) no agrega ruido, ya que deseamos una salida determinista.
\end{importantbox}

\subsection{Pseudocódigo del muestreo}

\begin{lstlisting}[caption={Algoritmo de muestreo DDPM}, label=lst:sample]
def ddpm_sample(model, T, alpha, alpha_bar, beta):
    x = torch.randn(shape)  # x_T ~ N(0, I)

    for t in range(T, 0, -1):
        epsilon_pred = model(x, t)

        # Compute mean
        mu = (1 / alpha[t].sqrt()) * (
            x - (beta[t] / (1 - alpha_bar[t]).sqrt()) * epsilon_pred
        )

        if t > 1:
            sigma = beta[t].sqrt()  # or tilde_beta[t].sqrt()
            z = torch.randn_like(x)
            x = mu + sigma * z
        else:
            x = mu  # No noise at final step

    return x
\end{lstlisting}

\begin{warningbox}{La velocidad de muestreo es el principal cuello de botella de DDPM}
El muestreo de DDPM requiere iterar $T$ pasos (usualmente $T = 1000$), y cada paso necesita una propagación hacia adelante completa de la red neuronal. Esto implica que generar una imagen requiere 1000 propagaciones hacia adelante, lo cual es aproximadamente tres órdenes de magnitud más lento que las únicas propagaciones hacia adelante de GAN. Trabajos posteriores como DDIM, DPM-Solver y otros se han centrado en reducir el número de pasos de muestreo, logrando una calidad cercana en tan solo 10--50 pasos.
\end{warningbox}

\subsection{Visualización de $\hat{x}_0$ durante el proceso de muestreo}

Al utilizar una red de predicción de ruido, en cada paso intermedio podemos utilizar la relación para calcular la estimación actual de $x_0$:

$$\hat{x}_0^{(t)} = \frac{x_t - \sqrt{1-\bar{\alpha}_t}\, \epsilon_\theta(x_t, t)}{\sqrt{\bar{\alpha}_t}}$$

En las etapas tempranas del proceso de muestreo ($t$ cercano a $T$), $\hat{x}_0^{(t)}$ es muy difuso; a medida que $t$ disminuye, la predicción se vuelve gradualmente más clara. Esta visualización ilustra intuitivamente la característica de «refinamiento iterativo» de los modelos de difusión.

\begin{knowledgebox}{Refinamiento iterativo y Progressive Refinement}
La predicción de $\hat{x}_0$ en cada paso puede entenderse como: la mejor conjetura sobre el resultado final bajo el nivel actual de información. A medida que avanza la eliminación de ruido, aumenta la cantidad de información disponible (y disminuye el ruido), por lo que la predicción se vuelve naturalmente más precisa. Esta característica de *progressive refinement* es muy útil en aplicaciones prácticas —por ejemplo, para guiar la dirección de generación en pasos intermedios (*classifier guidance* / *classifier-free guidance*)—.
\end{knowledgebox}

\subsection{Resumen de la sección}

El proceso de muestreo de DDPM es la «operación inversa» del proceso de entrenamiento: partiendo de puro ruido, se elimina el ruido gradualmente hasta obtener una imagen limpia. En cada paso se predice el ruido, se calcula la media y se agrega ruido aleatorio controlado. La lentitud del muestreo es la principal limitación, pero existen numerosos métodos de aceleración posteriores.

%% ========== Sección 8: Capacidad expresiva de los modelos de difusión ==========
